\chapter{System Design}

This chapter describes the complete design of the Smart Queue Management System. It covers the overall system architecture, the technology stack, the database design, the module design, the token status flow, and the API design. The system follows a three-tier client-server architecture with a clear separation between the presentation layer, the business logic layer, and the data layer.

\section{System Architecture}

The Smart Queue Management System is designed using a three-tier client-server architecture consisting of the following layers:

\begin{enumerate}
    \item \textbf{Presentation Layer (Frontend):} The user interface is built using ReactJS. It handles all user interactions, including token generation, queue status display, admin login, and the admin dashboard. The frontend runs as a Single Page Application (SPA) in the browser on port 3000.
    \item \textbf{Business Logic Layer (Backend):} The API is built using ASP.NET Core Web API (.NET 10). It processes all HTTP requests from the frontend, enforces business rules (such as sequential token numbering and status transitions), and communicates with the database. The backend runs on port 5000.
    \item \textbf{Data Layer (Database):} PostgreSQL is used as the relational database to store admin credentials and all queue token records persistently. Entity Framework Core (EF Core) is used as the Object-Relational Mapper (ORM) to interact with the database from the backend.
\end{enumerate}

The frontend communicates with the backend exclusively through HTTP REST API calls using the Axios library. The backend connects to the PostgreSQL database using EF Core with the Npgsql provider. CORS (Cross-Origin Resource Sharing) is enabled on the backend to allow the frontend (port 3000) to send requests to the backend (port 5000).

\section{Technology Stack}

Table~\ref{tab:techstack} lists all technologies and tools used in the development of the Smart Queue Management System along with their purpose.

\begin{table}[htbp]
    \centering
    \begin{tabular}{|c|c|c|}
        \hline
        \textbf{Component} & \textbf{Technology} & \textbf{Purpose} \\
        \hline
        Frontend Framework & ReactJS (v19) & Building the user interface \\
        \hline
        HTTP Client & Axios & Sending API requests from frontend \\
        \hline
        Client-Side Routing & React Router DOM (v7) & Multi-page navigation in SPA \\
        \hline
        Frontend Build Tool & Vite & Development server and bundler \\
        \hline
        Backend Framework & ASP.NET Core Web API & REST API server (.NET 10) \\
        \hline
        ORM & Entity Framework Core & Database access from backend \\
        \hline
        Database & PostgreSQL & Persistent relational data storage \\
        \hline
        DB Connector & Npgsql EF Core Provider & Connects .NET to PostgreSQL \\
        \hline
        Styling & Plain CSS & Frontend page styling \\
        \hline
    \end{tabular}
    \caption{Technology Stack of the Proposed System}
    \label{tab:techstack}
\end{table}

\section{Database Design}

The database is named \texttt{queue\_management} and consists of two tables: \textbf{Admins} and \textbf{Tokens}. The design is intentionally kept simple with minimal fields to serve the core requirements of the system.

\subsection{Admins Table}

The Admins table stores the login credentials for system administrators. It is used by the Admin Login API to authenticate admin users. Table~\ref{tab:admins} shows its structure.

\begin{table}[htbp]
    \centering
    \begin{tabular}{|c|c|c|}
        \hline
        \textbf{Field Name} & \textbf{Data Type} & \textbf{Description} \\
        \hline
        Id & SERIAL (Primary Key) & Auto-incremented unique identifier \\
        \hline
        Username & VARCHAR(100) NOT NULL & Admin login username \\
        \hline
        Password & VARCHAR(100) NOT NULL & Admin login password \\
        \hline
    \end{tabular}
    \caption{Admins Table Structure}
    \label{tab:admins}
\end{table}

\subsection{Tokens Table}

The Tokens table stores all queue token records generated by users. Each row represents one customer's token with its current lifecycle status. Table~\ref{tab:tokens} shows the complete structure.

\begin{table}[htbp]
    \centering
    \begin{tabular}{|c|c|c|}
        \hline
        \textbf{Field Name} & \textbf{Data Type} & \textbf{Description} \\
        \hline
        Id & SERIAL (Primary Key) & Auto-incremented unique identifier \\
        \hline
        TokenNumber & INT NOT NULL & Sequential queue number \\
        \hline
        CustomerName & VARCHAR(200) NOT NULL & Name of the customer \\
        \hline
        Status & VARCHAR(50) NOT NULL & Waiting / Active / Completed \\
        \hline
        CreatedAt & TIMESTAMP NOT NULL & Date and time token was created \\
        \hline
    \end{tabular}
    \caption{Tokens Table Structure}
    \label{tab:tokens}
\end{table}

\section{Module Design}

The system is divided into the following three functional modules:

\subsection{User Module}

The User Module is accessible to any visitor without login. It provides:
\begin{itemize}
    \item A form to enter the customer name and generate a digital queue token.
    \item Display of the assigned token number and initial status (Waiting).
    \item A live queue status page showing the currently active token and queue statistics.
\end{itemize}

\subsection{Admin Module}

The Admin Module is accessible only after authentication. It provides:
\begin{itemize}
    \item A login page with username and password verification against the database.
    \item A dashboard showing all tokens in a tabular format with their status.
    \item Controls to call the next waiting token, manually complete any token, or reset the entire queue.
\end{itemize}

\subsection{Dashboard Module}

The Dashboard is embedded in both the Queue Status page and the Admin Dashboard. It shows:
\begin{itemize}
    \item Total number of tokens currently in \textbf{Waiting} status.
    \item The \textbf{Active} token number being currently served.
    \item Total number of \textbf{Completed} tokens since last reset.
\end{itemize}

\section{Token Status Flow}

Each token in the system follows a defined lifecycle through three states:

\begin{center}
\textbf{Waiting} $\rightarrow$ \textbf{Active} $\rightarrow$ \textbf{Completed}
\end{center}

\begin{itemize}
    \item When a token is first generated, its status is set to \textbf{Waiting}.
    \item When the admin calls the next token, the system finds the next token with the lowest number and Waiting status, and sets it to \textbf{Active}. The previously Active token is automatically moved to \textbf{Completed}.
    \item A token can also be manually set to \textbf{Completed} directly from the admin dashboard without going through the Call Next flow.
    \item Only one token can be Active at any time.
\end{itemize}

\section{API Design}

The backend exposes a set of RESTful API endpoints organized under two route prefixes: \texttt{/api/token} for queue operations and \texttt{/api/admin} for authentication. Table~\ref{tab:api} provides the complete API reference.

\begin{table}[htbp]
    \centering
    \begin{tabular}{|c|c|c|}
        \hline
        \textbf{HTTP Method} & \textbf{Endpoint} & \textbf{Description} \\
        \hline
        POST & /api/token/generate & Generate a new queue token \\
        \hline
        GET & /api/token/all & Retrieve all tokens in order \\
        \hline
        GET & /api/token/current & Get the currently active token \\
        \hline
        GET & /api/token/stats & Get waiting, active, completed counts \\
        \hline
        POST & /api/token/next & Call the next waiting token \\
        \hline
        PUT & /api/token/complete/\{id\} & Mark a specific token as completed \\
        \hline
        DELETE & /api/token/reset & Delete all tokens (reset queue) \\
        \hline
        POST & /api/admin/login & Admin username and password login \\
        \hline
    \end{tabular}
    \caption{REST API Endpoints}
    \label{tab:api}
\end{table}

\section{Folder Structure}

The project follows a simple and clean folder structure that separates the frontend, backend, and database scripts into distinct directories. The backend follows the Controller-Model-DbContext pattern and the frontend follows a pages-based component structure. Table~\ref{tab:folders} summarizes the main directories.

\begin{table}[htbp]
    \centering
    \begin{tabular}{|c|c|}
        \hline
        \textbf{Directory} & \textbf{Contents} \\
        \hline
        backend/Controllers/ & AdminController.cs, TokenController.cs \\
        \hline
        backend/Models/ & Admin.cs, Token.cs \\
        \hline
        backend/Data/ & AppDbContext.cs \\
        \hline
        frontend/src/pages/ & Home, GenerateToken, QueueStatus, AdminLogin, AdminDashboard \\
        \hline
        frontend/src/services/ & api.js (all Axios API call functions) \\
        \hline
        database/ & setup.sql \\
        \hline
    \end{tabular}
    \caption{Project Folder Structure}
    \label{tab:folders}
\end{table}
